\chapter{Problemas e gráficos}\label{ch:g5:data}

O último capítulo do ano põe tudo para funcionar: problemas de várias
etapas com dinheiro, quantidades e durações, e os gráficos que
apresentam a informação num relance. A ideia-estrela é ``por
unidade'' --- descobrir o preço de um só item --- que vai crescer e
virar proporcionalidade no \cref{ch:g6:propdata}.

\section{Problemas de várias etapas}

\begin{method}[Domar um problema comprido]\label{met:g5:data:steps}
\begin{enumerate}
  \item Leia até o fim; diga qual é a pergunta final;
  \item planeje de trás para frente: para responder a ela, do que
        preciso primeiro?
  \item resolva as etapas em ordem, escrevendo uma operação e uma
        frase curta por etapa;
  \item confira a resposta final com o bom senso e com uma
        estimativa.
\end{enumerate}
\end{method}

\begin{example}\label{ex:g5:data:multistep}
``Uma escola encomenda $4$ caixas de $25$ livros a $6$ cada um e paga
$180$ de entrega por cima. Qual é o custo total?''
\begin{enumerate}
  \item Livros: $4 \times 25 = 100$ livros.
  \item Custo dos livros: $100 \times 6 = 600$.
  \item Total: $600 + 180 = 780$.
\end{enumerate}
Conferência por estimativa: uns cem livros a $6$ dão uns $600$, mais
uns $200$: uns $800$ --- a resposta $780$ é plausível.
\end{example}

\begin{example}[Por unidade]\label{ex:g5:data:perone}
``$5$ canecas iguais custam $12.50$; quanto custam $8$ canecas?''
Descubra primeiro o preço de \emph{uma} caneca:
\[
12.50 \div 5 = 2.50,
\qquad\text{e depois}\qquad
8 \times 2.50 = 20 .
\]
Pela unidade, tudo fica ao alcance --- esse padrão de duas etapas
resolve toda uma família de problemas (receitas, velocidades,
preços) e é a semente do \cref{ch:g6:propdata}.
\end{example}

\begin{example}[Troco e orçamento]\label{ex:g5:data:budget}
Zoe tem $30$. Ela compra um livro a $12.40$ e duas canetas a $2.30$
cada. Ela ainda consegue comprar um jogo de tabuleiro de $13$?
\begin{enumerate}
  \item Canetas: $2 \times 2.30 = 4.60$.
  \item Gasto até aqui: $12.40 + 4.60 = 17$.
  \item Sobra: $30 - 17 = 13$ --- exatamente o suficiente para o
        jogo, sem sobrar nada.
\end{enumerate}
\end{example}

\section{Ler gráficos}

\begin{method}[Ler qualquer gráfico]\label{met:g5:data:read}
\begin{enumerate}
  \item Leia o título: o que está sendo contado?
  \item Leia os eixos ou a legenda: o que vale uma graduação, uma
        barra, um símbolo?
  \item Só então responda às perguntas --- apontar o gráfico com uma
        régua ajuda a ler as alturas com exatidão.
\end{enumerate}
\end{method}

\begin{example}\label{ex:g5:data:chart}
Chuva numa cidade, mês a mês:

\begin{omfigure}
\begin{tikzpicture}
\begin{axis}[
  width=8.8cm, height=5.0cm,
  ybar, bar width=13pt,
  symbolic x coords={Jan,Fev,Mar,Abr,Mai,Jun},
  xtick=data,
  ymin=0, ymax=90,
  ytick={20,40,60,80},
  ylabel={chuva (mm)},
  tick label style={font=\small},
  label style={font=\small}]
  \addplot[fill=omDef!30, draw=omDef] coordinates
    {(Jan,60) (Fev,45) (Mar,55) (Abr,70) (Mai,80) (Jun,30)};
\end{axis}
\end{tikzpicture}

{\small Um gráfico de barras da chuva. Uma graduação $= 20$ mm.}
\end{omfigure}

Leituras: o mês mais chuvoso é maio ($80$ mm); em junho choveu a
\omterm{def:g3:division:half}{metade} do que em abril ($30$ contra $70$)? Não --- a \omterm{def:g3:division:half}{metade} de $70$ é
$35$, então \emph{não chega} à \omterm{def:g3:division:half}{metade}: ler com precisão importa!
Total dos seis meses:
$60 + 45 + 55 + 70 + 80 + 30 = 340$ mm.
\end{example}

\begin{example}[Uma grandeza que muda]\label{ex:g5:data:line}
A altura de uma planta, medida toda semana:

\begin{omfigure}
\begin{tikzpicture}
\begin{axis}[omaxis,
  width=8.4cm, height=5.0cm,
  xmin=0, xmax=6.5, ymin=0, ymax=24,
  xlabel={semana}, ylabel={altura (cm)},
  xtick={1,2,3,4,5,6}, ytick={5,10,15,20}]
  \addplot[omDef, thick, mark=*, mark size=1.6pt] coordinates
    {(1,3) (2,6) (3,10) (4,15) (5,19) (6,21)};
\end{axis}
\end{tikzpicture}

{\small Unir os pontos de medida mostra o \emph{crescimento}: quanto
mais inclinado o segmento, mais depressa a planta cresceu naquela
semana.}
\end{omfigure}

A planta cresceu mais depressa entre as semanas $3$ e $4$ ($+5$ cm) e
desacelerou no fim ($+2$ cm na semana $6$).
\end{example}

\section{Exercícios}

\begin{exercise}[$\star$]\label{exo:g5:data:1}
Resolva por etapas: ``Um clube aluga um ônibus por $240$ e compra
$32$ ingressos de museu a $7$ cada. Qual é o custo total do
passeio?''
\end{exercise}

\begin{exercise}[$\star$]\label{exo:g5:data:2}
Resolva com o ``por unidade'': ``$3$ kg de maçãs custam $7.50$.
Quanto custam $5$ kg?''
\end{exercise}

\begin{exercise}[$\star$]\label{exo:g5:data:3}
Resolva: ``Seis amigos dividem igualmente o custo de $87$ de um
piquenique. Quanto cada um paga?''
\end{exercise}

\begin{exercise}[$\star$]\label{exo:g5:data:4}
Samuel tem $25$. Ele compra uma camiseta a $13.90$ e um boné a
$8.50$. Quanto dinheiro sobra para ele?
\end{exercise}

\begin{exercise}[$\star$]\label{exo:g5:data:5}
Usando o gráfico de chuva do \cref{ex:g5:data:chart}: em que meses
choveu mais de $50$ mm? Quanta chuva a mais caiu em maio do que em
junho?
\end{exercise}

\begin{exercise}[$\star$]\label{exo:g5:data:6}
Usando o gráfico da planta do \cref{ex:g5:data:line}: que altura a
planta tinha na semana $2$? Entre que semanas ela cresceu exatamente
$4$ cm?
\end{exercise}

\begin{exercise}[$\star$]\label{exo:g5:data:7}
Faça um gráfico de barras para o número de gols de um time:
setembro $4$, outubro $7$, novembro $3$, dezembro $6$. Escolha a
graduação e dê o total.
\end{exercise}

\begin{exercise}[$\star$]\label{exo:g5:data:8}
Um pacote de $8$ iogurtes custa $3.20$; vendido sozinho, um iogurte
custa $0.50$. Calcule o preço ``por unidade'' no pacote e quanto se
economiza por iogurte comprando o pacote.
\end{exercise}

\begin{exercise}[$\star\star$]\label{exo:g5:data:9}
``Um cinema vendeu $148$ ingressos a $9.50$ no sábado e $95$
ingressos ao mesmo preço no domingo. Quanto dinheiro a mais o sábado
rendeu em relação ao domingo?'' Resolva de duas maneiras: calculando
a arrecadação dos dois dias, ou calculando com a \emph{\omterm{ex:g1:subtraction:difference}{diferença}} de
ingressos --- e confira que as respostas coincidem.
\end{exercise}

\begin{exercise}[$\star\star$]\label{exo:g5:data:10}
Uma receita para $4$ pessoas pede $300$ g de arroz. Aline cozinha
para $10$ pessoas. De quanto arroz ela precisa? (Primeiro por pessoa
--- decimais são permitidos.)
\end{exercise}

\begin{exercise}[$\star\star$]\label{exo:g5:data:11}
Invente um problema de três etapas com dinheiro cuja resposta final
seja $5.50$, escreva a solução completa no estilo do
\cref{met:g5:data:steps} e teste-o com um colega.
\end{exercise}
